\documentclass[10pt]{article}

\usepackage[T1]{fontenc}
\usepackage[margin=0.72in]{geometry}
\usepackage{amsmath}
\usepackage{array}
\usepackage{booktabs}
\usepackage{enumitem}
\usepackage{xcolor}
\usepackage{hyperref}
\usepackage{tabularx}

\hypersetup{hidelinks}
\setlength{\parindent}{0pt}
\setlength{\parskip}{0.45em}
\setlist{nosep,leftmargin=1.4em}
\renewcommand{\arraystretch}{1.13}

\definecolor{implemented}{HTML}{26734D}
\definecolor{target}{HTML}{1D4ED8}
\definecolor{unproven}{HTML}{9A3412}
\newcommand{\Implemented}{\textcolor{implemented}{\textbf{IMPLEMENTED}}}
\newcommand{\Target}{\textcolor{target}{\textbf{TARGET}}}
\newcommand{\Unproven}{\textcolor{unproven}{\textbf{UNPROVEN}}}
\newcommand{\code}[1]{\texttt{#1}}

\title{Claude VM Performance and Non-Interference Contract}
\author{Lee Daghlar Ostadi}
\date{Contract version 0.1 --- 2026-09-01}

\begin{document}
\maketitle
\vspace{-1.8em}

\section{Purpose and status vocabulary}

This document defines the performance-only contract for keeping one Claude Linux
VM warm for the lifetime of its OpenCode controller.  It changes neither the
Qwen model nor any access capability.  It does not authorize a new device,
filesystem, network, USB, signing, execution, or policy surface.

The following labels are normative:

\begin{center}
\begin{tabularx}{\linewidth}{@{}>{\bfseries}l X@{}}
\toprule
\Implemented & Present in the inspected implementation or supported by a recorded receipt. \\
\Target & Required behavior that must pass the acceptance gates in Section~\ref{sec:gates}. \\
\Unproven & A claim for which the current evidence is insufficient.  It must not be reported as implemented. \\
\bottomrule
\end{tabularx}
\end{center}

The contract distinguishes configuration equality from performance equality.
Byte-for-byte configuration stability is mandatory.  Performance is accepted
only through repeated measurements under the same host conditions.

\section{Entities, clocks, and baseline}

Let \(C\) be the OpenCode-launched \code{claude\_vm} MCP controller process,
\(V\) its single \code{ClaudeVZRunner} child and Apple
\code{VZVirtualMachine}, and \(Q\) the q3.8 Qwen model selected when OpenCode
starts.  Let \(L_C=[t_C^{start},t_C^{exit}]\) be the controller lifetime.
Model turns, final responses, compaction, and Ollama load or unload events are
not controller-exit events.

Guest-ready latency is
\[
  T_{ready}=t(\code{guest\_ready})-t(\code{starting}).
\]
The recorded current reference is
\[
  B_{ready}=5.075\ \mathrm{s}.
\]
This value is an \Implemented{} baseline reference, not permission to replace
raw timestamps.  Every qualification run must retain the corresponding
\code{starting}, \code{vz\_started}, and \code{guest\_ready} timestamps.
Reproducibility of the reference under the final persistent design remains
\Unproven{} until the gates in Section~\ref{sec:gates} pass.

The current compound execution path starts and stops the VM around an operation.
Its guest-ready and cleanup receipts are \Implemented.  Session-long residency
and coexistence with unchanged Qwen are \Target{} behavior.

The process-lifetime harness reached \code{guest\_ready} in 3.551~s, retained
one runner and lease across 20 operations, and recorded a 0.274~s median,
0.373~s p95, and 0.520~s maximum operation latency.  Controller close removed
the runner, disk handles, and lease in 0.744~s.  This is \Implemented{}
harness-level evidence for single initialization, warm residency, and
controller cleanup.  Actual OpenCode-window closure remains a separate gate.

A separate unchanged-Qwen coexistence probe failed the non-interference gate.
The no-VM control used a 16.687~s model-load phase.  With the 4~GiB VM warm,
Qwen remained the same model at context 262144 and used a 14.721~s load phase,
but system free-memory headroom fell to 12\%.  Total responses are not compared
because the generated token counts differed.  The memory result fails this
contract regardless of the faster single load sample.  Consequently, local
process-lifetime mode remains disabled in the OpenCode registration pending a
design that passes every gate.

\section{Process-owned lifetime}

\subsection*{Ownership invariant}

The controller \(C\) is the sole lifetime owner of \(V\).  Exactly one VM lease,
one runner process, and one VM instance may exist for \(L_C\).  Model output and
model tool selection do not own this lifetime.

\begin{enumerate}
  \item \Target{} Controller initialization starts \(V\) exactly once and waits
        for application-level \code{guest\_ready}.  The first model operation
        may observe either ``initializing'' or ``ready'', but may not create a
        second VM.
  \item \Target{} Once ready, \(V\) remains running throughout \(L_C\).
        Completion of a prompt, a final model response, an idle interval, Qwen
        compaction, or Qwen unload is not a VM stop condition.
  \item \Target{} Guest operations reuse the existing runner, lease, vsock
        connection, handshake, and ready state.  An operation may not call the
        current start--execute--stop cycle internally.
  \item \Target{} Only controller exit initiates normal teardown.  OpenCode
        window closure must cause controller stdin EOF, \code{SIGINT}, or
        \code{SIGTERM}; the controller then performs bounded stop, confirms no
        clone-disk handles, and releases its lease without another model turn.
  \item \Implemented{} Signal and stdin-EOF cleanup handlers already exist.
        The exact mapping from one OpenCode prompt-window closure to controller
        exit is \Unproven{} until observed end to end.
\end{enumerate}

An unexpected runner failure may end \(V\), but must be recorded as a failed
lifetime rather than silently replaced during a model call.  A deliberate
restart, if later specified, is a new measured lifetime and is outside this
contract version.

\section{Warm reuse and non-interference}

For \(N\geq20\) allowed guest operations in one controller lifetime, the event
stream must contain one \code{starting}, one \code{vz\_started}, one
\code{guest\_ready}, and no intermediate \code{stopping}, \code{stopped}, or
runner \code{process\_exit} event.  The runner PID, machine identifier, lease
owner, and VM identity remain constant.  There is one bounded stop sequence only
after controller exit.

Warm reuse removes repeated \(T_{ready}\) from per-operation latency.  It does
not relax the existing operation timeout, argument, output, correlation, or
serialization bounds.  Large guest output remains external to Qwen except for
the already bounded receipt.  No performance claim may be based on adding a
tool schema, instruction block, or guest transcript to Qwen's initial context.

The following Qwen launch snapshot is immutable for the experiment:

\begin{itemize}
  \item provider, model tag, resolved model manifest, and model identity;
  \item configured context capacity, output limit, and model parameters;
  \item Ollama runtime settings, including KV type, Flash Attention,
        parallelism, loaded-model limit, and keep-alive;
  \item OpenCode model selection, permissions, MCP registry, instruction files,
        skill roots, and effective configuration.
\end{itemize}

Before and after each qualification run, these fields are serialized in a
stable order and hashed as \(H_Q\).  Acceptance requires
\(H_Q^{before}=H_Q^{after}\).  Reducing Qwen's context, changing its model,
changing its prompt/tool surface, or tuning its runtime invalidates the run.
The persistent VM must not issue an Ollama unload, reload, or configuration
operation.  Qwen may complete its normal lifecycle independently, but it may not
stop or unload \(V\).

For a fixed probe prompt, input-token count must remain identical before and
during VM residency.  Median time to first token and median generation rate are
measured in paired trials; the persistent VM passes only if time to first token
regresses by no more than 5\% and generation rate regresses by no more than 5\%.
No model eviction, context reset, out-of-memory event, or critical host-memory
pressure is permitted in a passing trial.

\section{Frozen capability surface}

The performance implementation must preserve the following current surface.
Each row is fingerprinted before and after qualification; a mismatch fails the
run rather than being treated as a performance improvement.

\begin{center}
\small
\begin{tabularx}{\linewidth}{@{}>{\raggedright\arraybackslash}p{0.19\linewidth}
                                  >{\raggedright\arraybackslash}p{0.18\linewidth} X@{}}
\toprule
Surface & Status & Frozen value \\
\midrule
Policy and limits & \Implemented & Current \code{policy.json}, including 4 GiB, 4 vCPU, 30 s startup timeout, argument limits, 32 KiB output ceiling, and serialized operation policy. \\
Network & \Implemented & \code{isolated} blackhole Virtio network; no NAT, bridge, forwarding, host reachability, DNS service, or external egress. \\
Host filesystem & \Implemented & One read-only \code{claudeshared} VirtioFS directory; no additional host path. \\
VM devices & \Implemented & Current storage, entropy, memory-balloon, vsock, network, console, and VirtioFS topology. \\
USB & \Implemented & No guest USB device is configured by the local Swift runner.  This contract neither adds nor removes USB access.  Any broader USB claim is \Unproven{} and outside this performance change. \\
Signing & \Implemented & The local runner's current ad-hoc signature and current \code{com.apple.security.virtualization} entitlement.  No signature or entitlement is inherited from the original Claude application. \\
Guest execution & \Implemented & Exact current allowlist only: \code{uname} with \code{[]}, \code{-a}, \code{-m}, \code{-r}, or \code{-s}; \code{id} with \code{[]}, \code{-u}, or \code{-g}; \code{python3 --version}; and \code{git --version}. \\
\bottomrule
\end{tabularx}
\end{center}

The target changes when the VM starts and stops, not what it can reach or run.
Any access-surface change requires a separate contract and cannot be included in
the measurements defined here.

\section{Acceptance gates}\label{sec:gates}

All gates are required.  A result set records wall-clock timestamps, process
identifiers, VM events, Qwen snapshot hashes, model measurements, capability
fingerprints, and final cleanup state.

\begin{enumerate}[label=G\arabic*.]
  \item \textbf{Cold-ready regression.} Run at least five alternating control
        and candidate cold starts under the same host state.  Candidate median
        \(T_{ready}\) must be no greater than
        \(\max(5.075\mathrm{s},\operatorname{median}(control))+0.250\mathrm{s}\),
        and every run must remain below the unchanged 30 s startup timeout.

  \item \textbf{Single initialization.} Across at least 20 operations in one
        controller lifetime, event counts must be exactly one start, one
        \code{vz\_started}, and one \code{guest\_ready}; no operation may incur
        another boot or readiness handshake.

  \item \textbf{Warm residency.} The same runner PID, VM identity, lease owner,
        and ready vsock connection must remain live between all operations.
        Aggregate session boot cost is therefore one measured \(T_{ready}\),
        not \(N\times T_{ready}\).

  \item \textbf{Qwen invariance.} \(H_Q^{before}=H_Q^{after}\); the fixed probe
        has identical input-token count; no Qwen model reload, eviction, context
        reset, or configuration mutation occurs.

  \item \textbf{Qwen performance.} In at least ten paired fixed-prompt trials,
        candidate median time to first token is at most 105\% of control, and
        candidate median generation rate is at least 95\% of control.  Any
        out-of-memory or critical-pressure event fails the gate.  After Qwen is
        fully loaded, host free-memory headroom must remain at least 35\% and
        the paired interval must add zero swapouts.

  \item \textbf{Capability invariance.} Policy, device topology, network mode,
        USB surface, signing identity and entitlements, host share, and allowlist
        fingerprints must match their pre-change values exactly.

  \item \textbf{Controller-owned shutdown.} Closing the tested OpenCode prompt
        window must terminate its controller and initiate the only normal VM
        stop.  Within 60 s, the runner and Apple VM XPC process are absent, both
        clone disks have no open handle, and the controller lease is absent.
        No model tool call may be required for cleanup.

  \item \textbf{Receipt completeness.} The final receipt records all preceding
        gates and labels failures honestly.  A successful warm operation alone
        does not establish process-owned lifetime, non-interference, or shutdown.
\end{enumerate}

\section{Current-to-target summary}

\begin{center}
\begin{tabularx}{\linewidth}{@{}l X@{}}
\toprule
Status & Claim \\
\midrule
\Implemented & Apple Virtualization start, application-level guest readiness,
bounded allowlisted execution, serialized controller ownership, bounded stop,
and cleanup receipts exist.  The current readiness reference is 5.075 s. \\
\Target & One controller starts one VM, reuses it warm without model-directed
stop, leaves Qwen and every capability fingerprint unchanged, and stops the VM
only when that OpenCode controller exits. \\
\Unproven & Full-lifetime Qwen/VM coexistence without measurable regression;
the prompt-window-to-controller-exit mapping; repeatability of the 5.075 s
reference under the target lifecycle; and any guest USB capability.  The first
local coexistence candidate failed because free-memory headroom reached 12\%. \\
\bottomrule
\end{tabularx}
\end{center}

This contract is satisfied only by the complete evidence set.  It does not
permit substituting a smaller model, shorter context, changed runtime, expanded
access surface, or per-call VM cycle for the target behavior.

\end{document}
